%% ma: L10-B13-0092
%% lop: 10
%% chuong: 5
%% bai: 13
%% dang: 10.13.10
%% loai: TN4
%% muc: NB
%% dap_an: "B"
%% nguon: "chuyen-de-cac-so-dac-trung-cua-mau-so-lieu-khong-ghep-nhom-toan-10-chuong-trinh-moi.pdf (GV. Phan Nhật Linh) (P03), Bài 02 Dạng 2 Phần I Câu 19"
%% kien_thuc: "Lớp 10 Bài 13 (trung vị, tứ phân vị)"
%% trang_thai: chinh-thuc
%% ghi_chu: "Tự duyệt theo yêu cầu giáo viên 22/10/2026."
\begin{ex}
Số học sinh của các lớp khối 12 trong một trường THPT được thống kê như sau: $35$; $40$; $37$; $46$; $38$; $42$; $43$; $36$. Tứ phân vị thứ hai của mẫu số liệu là
\choice{$40$}{\True $39$}{$36$}{$42$}
\loigiai{Sắp xếp: $35$; $36$; $37$; $38$; $40$; $42$; $43$; $46$. $Q_2=\dfrac{38+40}2=39$. Chọn B.}
\end{ex}
